\documentclass[UTF8,11pt]{ctexart}

\usepackage[a4paper,margin=1in]{geometry}
\usepackage{amsmath,amssymb,bm,mathtools}
\usepackage{booktabs}
\usepackage{hyperref}
\hypersetup{colorlinks=true,linkcolor=blue,urlcolor=blue,citecolor=blue}

\title{MAGT 第 3 部分证明讲解与随机插值一步生成的可移植性分析}
\author{}
\date{基于 arXiv:2602.19600v1}

\newcommand{\E}{\mathbb{E}}
\newcommand{\R}{\mathbb{R}}
\newcommand{\cH}{\mathcal{H}}
\newcommand{\cU}{\mathcal{U}}
\newcommand{\cM}{\mathcal{M}}
\newcommand{\cN}{\mathcal{N}}
\newcommand{\dd}{\,\mathrm{d}}
\newcommand{\KL}{\mathrm{KL}}
\newcommand{\op}{\mathrm{op}}
\newcommand{\supp}{\mathrm{supp}}
\newcommand{\Ent}{\mathrm{Ent}}
\newcommand{\Cov}{\mathrm{Cov}}

\begin{document}

\maketitle

\section{阅读目标}

论文第 3 部分的目标是证明：如果一个低维到高维的非可逆传输映射
\[
h:\R^d\to\R^D,\qquad d\ll D
\]
在一个固定 Gaussian smoothing level \(t\) 上学到了足够准确的 smoothed score，那么它的一步采样
\[
\widetilde Y_0=\hat h_\lambda(U),\qquad U\sim\pi
\]
在未加噪的原始空间中也能以 \(W_2\) 距离接近真实数据分布。

这一部分的证明不是单一技术，而是四层闭环：
\[
\text{fixed-level score risk}
\Longrightarrow
\text{Fisher mismatch}
\Longrightarrow
W_2(P_{Y_0},P_{\widetilde Y_0})
\Longrightarrow
\text{finite-sample rate}.
\]

更具体地，第 3 部分由三个主要定理组成：
\begin{enumerate}
\item \textbf{Single-level pull-back bound}：把固定噪声层级的 score mismatch 拉回到未加噪分布的 \(W_2\) 误差。
\item \textbf{Score-matching excess-risk bound}：用 bracketing entropy 控制经验风险最小化器的 excess risk。
\item \textbf{Generation fidelity theorem}：合并 pull-back、经验过程、神经网络逼近和 anchor 近似误差，得到有限样本生成误差界。
\end{enumerate}

\section{模型设定与几何假设}

论文假设真实数据由一个低维 latent 经过未知真传输 \(h^*\) 生成：
\[
Y_0=h^*(U),\qquad U\sim\pi,\qquad h^*:\cU\subset\R^{d^*}\to\R^D.
\]
真实分布支撑在低维流形
\[
\cM^*=h^*(\supp\pi)
\]
上或附近。学习目标是用 \(\hat h_\lambda\) 生成
\[
\widetilde Y_0=\hat h_\lambda(U).
\]

核心几何假设有三类。

\subsection{正则传输类}

要求 \(h^*\in C^{\eta+1}(\cU,B)\)，并且候选函数 \(h\in\cH_{\mathrm{reg}}\) 满足 Jacobian 满秩和条件数控制：
\[
\mathrm{rank}\,J_h(u)=d,\qquad
\sigma_{\min}(J_h(u))\ge m,\qquad
\sigma_{\max}(J_h(u))\le M.
\]
同时 \(J_h\) 是 \(C^{0,\gamma}\) 的，其中
\[
\gamma=\min(1,\eta).
\]

这组假设的作用是：保证 \(h(\cU)\) 真的是一个规则的 \(d\)-维嵌入流形，并且局部坐标图不会塌缩。

\subsection{正 reach}

要求流形 \(\cM\) 的 reach 至少为 \(\rho_\cM>0\)。reach 可以理解为“管状邻域内最近点投影唯一”的半径。它保证当噪声不太大时，点 \(y/\alpha_t\) 可以唯一投影回流形附近。

这对附录中的 Hessian bound 至关重要，因为证明会把 \(y/\alpha_t\) 投影到流形上，然后在投影点附近做局部坐标展开。

\subsection{latent prior 的 LSI}

要求 \(\pi\) 满足 log-Sobolev inequality：
\[
\Ent_\pi(f^2)
\le
\frac{2}{C_{\mathrm{LSI}}(\pi)}
\int \|\nabla f\|^2\,\dd\pi.
\]

这个假设的作用不是为了优化算法，而是为了把 Fisher divergence 转换成 \(W_2\) 距离。附录证明中先把 \(\pi\otimes \cN(0,I_D)\) 的 LSI 通过 Lipschitz 映射
\[
F(u,z)=\alpha_t h(u)+\sigma_t z
\]
推到 smoothed law 上，从而得到
\[
C_{\mathrm{LSI}}(q_t)
\ge
\frac{\min\{C_{\mathrm{LSI}}(\pi),1\}}
{\alpha_t^2M^2+\sigma_t^2}.
\]

\section{固定层级 smoothing 与 denoiser mismatch}

第 3 节使用 VP smoothing：
\[
Y_t=\alpha_tY_0+\sigma_tZ,
\qquad
\widetilde Y_t=\alpha_t\widetilde Y_0+\sigma_tZ,
\qquad
Z\sim\cN(0,I_D),
\]
其中
\[
\alpha_t=\sqrt{1-t},
\qquad
\sigma_t^2=t.
\]

令 \(p_t\) 和 \(\widetilde p_t\) 分别为 \(Y_t\) 和 \(\widetilde Y_t\) 的密度。因为加了 Gaussian noise，哪怕 \(Y_0\) 支撑在低维流形上，\(p_t\) 也在 ambient space 中有良好密度。

论文定义 denoiser mismatch：
\[
E_{\mathrm{MAG}}(t)
:=
\E_{Y\sim p_t}
\left\|
m_{p,t}(Y)-m_{\widetilde p,t}(Y)
\right\|_2^2,
\]
其中
\[
m_{p,t}(y)=\E[Y_0\mid Y_t=y],
\qquad
m_{\widetilde p,t}(y)=\E[\widetilde Y_0\mid \widetilde Y_t=y].
\]

Tweedie 公式给出：
\[
m_{p,t}(y)
=
\frac{y+t\nabla\log p_t(y)}{\alpha_t},
\qquad
m_{\widetilde p,t}(y)
=
\frac{y+t\nabla\log \widetilde p_t(y)}{\alpha_t}.
\]
因此
\[
E_{\mathrm{MAG}}(t)
=
\frac{t^2}{1-t}
\E_{Y\sim p_t}
\left\|
\nabla\log p_t(Y)-\nabla\log\widetilde p_t(Y)
\right\|_2^2.
\]

也就是说，denoiser mismatch 与 Fisher divergence 等价：
\[
E_{\mathrm{MAG}}(t)
=
\frac{t^2}{1-t}\,
\mathcal J(p_t\|\widetilde p_t).
\]

\section{定理一：Single-level pull-back bound}

定理一的结论是
\[
W_2(P_{Y_0},P_{\widetilde Y_0})
\le
C_{\mathrm{PB}}(t)\sqrt{E_{\mathrm{MAG}}(t)}.
\]
其中
\[
C_{\mathrm{PB}}(t)
=
\frac{\sqrt{1-t}}{t}
\left(
\Phi(t)\bar C_{\mathrm{LSI}}(t)+\Psi(t)
\right).
\]

这一定理是整篇第 3 部分的核心。它说明只要在一个固定噪声层级 \(t\) 上控制 score 或 denoiser mismatch，就能控制未加噪的一步生成分布。

\subsection{证明第 1 步：引入 VP probability flow}

附录把 \(t\) 改写为标准时间参数 \(s\)：
\[
t=\sigma^2(s)=1-e^{-s},
\qquad
s=\Gamma(t)=-\log(1-t).
\]
记 \(T=\Gamma(t)\)。在 VP probability flow 中，密度 \(p_s\) 的速度场为
\[
v_s(x)
=
-\frac12x-\nabla\log p_s(x)
\]
在常数 schedule 下的形式。

对真实 smoothed density \(p_s\) 和模型 smoothed density \(\widetilde p_s\)，附录使用一个变量系数 pull-back 引理：
\[
W_2(p_0,q_0)
\le
e^{\int_0^T L}
W_2(p_T,q_T)
+
e^{\int_0^T L}
\int_0^T
e^{\int_u^T(K-L)}
\dd u\,
\sqrt{J_T}.
\]
这里
\[
J_T
=
\int
\left\|
\nabla\log p_T-\nabla\log q_T
\right\|^2q_T.
\]

直觉是：如果两个分布在时间 \(T\) 已经接近，并且 score gap 沿着反向 probability flow 不会爆炸，那么可以把时间 \(T\) 的误差拉回到时间 \(0\)。

\subsection{证明第 2 步：用 LSI 控制平滑层的 \(W_2\)}

pull-back 引理右侧第一项含有
\[
W_2(p_t,\widetilde p_t).
\]
附录用 LSI \(\Rightarrow\) Talagrand \(\Rightarrow\) Fisher 链式不等式：
\[
W_2^2(q,p)
\le
\frac{2}{\rho}\KL(q\|p)
\le
\frac{1}{\rho^2}\mathcal J(q\|p).
\]
因此
\[
W_2(p_t,\widetilde p_t)
\le
\bar C_{\mathrm{LSI}}(t)\sqrt{J_T}.
\]

这里
\[
\bar C_{\mathrm{LSI}}(t)
\le
\frac{(1-t)M^2+t}
{\min\{C_{\mathrm{LSI}}(\pi),1\}}.
\]

\subsection{证明第 3 步：Hessian bound 给出 \(L(s)\) 和 \(K(s)\)}

pull-back 引理需要两个控制量：
\[
\|\nabla v_s\|_{\op}\le L(s),
\qquad
\frac{\dd}{\dd s}J_s\le 2K(s)J_s.
\]

附录用 Hessian bound 建立
\[
\|\nabla^2\log p_s\|_{\op}
\vee
\|\nabla^2\log \widetilde p_s\|_{\op}
\le
L_*(s),
\]
其中
\[
L_*(s)
=
C_T^{(\gamma)}
\frac{\sigma_s^{\gamma-2}}{\alpha_s^\gamma}
+
\frac{(C_S^{(\gamma)})^2}{1-\theta_t}
\frac{\sigma_s^{2\gamma-2}}{\alpha_s^{2\gamma}}.
\]

由于
\[
v_s(x)=-\frac12x-\nabla\log\rho_s(x),
\]
可得
\[
L(s)=\frac12+L_*(s).
\]
score-gap growth lemma 进一步给出
\[
K(s)=\frac12+4L_*(s).
\]

\subsection{证明第 4 步：几何 Hessian bound 的来源}

Hessian bound 是最几何的一步。它从 Tweedie Hessian 公式出发：
\[
\nabla^2\log p_s(y)
=
\frac{\alpha_s^2}{\sigma_s^4}
\Cov(X\mid Y_s=y)
-
\frac1{\sigma_s^2}I_D.
\]

然后把 \(y/\alpha_s\) 投影到流形上：
\[
x_0=\Pi_{\cM}(y/\alpha_s).
\]
reach 条件保证投影唯一。再在 \(x_0\) 附近做局部 chart：
\[
x(\xi)
=
x_0+J\xi+R(\xi),
\qquad
\|R(\xi)\|\le C\|\xi\|^{1+\gamma}.
\]

条件后验在切向方向的协方差近似是
\[
\Cov_T
=
\frac{\sigma_s^2}{\alpha_s^2}\Pi_T
+
O\left(
\frac{\sigma_s^{2+\gamma}}{\alpha_s^{2+\gamma}}
\right),
\]
法向方向的协方差更小：
\[
\Cov_N
=
O\left(
\frac{\sigma_s^{2+2\gamma}}{\alpha_s^{2+2\gamma}}
\right).
\]

代回 Hessian 公式后，切向上的主项抵消，而法向上保留强负曲率：
\[
\Pi_NH_s\Pi_N
\preceq
-
\left(
\sigma_s^{-2}
-
C\frac{\sigma_s^{2\gamma-2}}{\alpha_s^{2\gamma}}
\right)\Pi_N.
\]
这就是为什么证明需要 \(t\) 落在 tube regime 中：噪声不能大到离开流形的唯一投影邻域。

\subsection{证明第 5 步：把 \(\sqrt{J_T}\) 换成 \(\sqrt{E_{\mathrm{MAG}}(t)}\)}

最后用
\[
E_{\mathrm{MAG}}(t)
=
\frac{t^2}{1-t}J_T
\]
得到
\[
\sqrt{J_T}
=
\frac{\sqrt{1-t}}{t}
\sqrt{E_{\mathrm{MAG}}(t)}.
\]
代回 pull-back 结果，即得定理一。

\section{定理二：Score-matching excess-risk bound}

第 3 节定义理想无限 anchor 损失：
\[
\ell(y_t,y_0;h)
=
\left\|
s_t(y_t;h)
-
\nabla_{y_t}\log p(y_t\mid y_0)
\right\|_2^2,
\]
其中
\[
s_t(\cdot;h)=\nabla\log p_t^h(\cdot).
\]

总体风险为
\[
R(h)=\E[\ell(Y_t,Y_0;h)].
\]
excess risk 定义为
\[
\rho^2(h^*,h)
=
R(h)-R(h^*).
\]

\subsection{正交分解}

训练 target 是 conditional score：
\[
\nabla_{y_t}\log p(Y_t\mid Y_0).
\]
它满足无偏性：
\[
\E[
\nabla_{y_t}\log p(Y_t\mid Y_0)
\mid Y_t
]
=
\nabla\log p_t(Y_t).
\]
于是得到正交分解：
\[
R(h)
=
R(h^*)
+
\E_{Y\sim p_t}
\left\|
s_t(Y;h)-s_t(Y;h^*)
\right\|_2^2.
\]
所以
\[
\rho^2(h^*,h)
=
\mathcal J(p_t\|p_t^h).
\]

这一点非常关键：它使得经验 score-matching 风险真正校准到固定层级 Fisher divergence。

\subsection{经验过程证明结构}

实际训练用有限 anchor 损失 \(\ell_K\)。假设它与理想损失的偏差满足
\[
\Delta_K
:=
\sup_{h,x}|\ell_K(x;h)-\ell(x;h)|
\le
\varepsilon^2/8.
\]

令 \(h^\#\) 是函数类 \(\cH\) 中的 population best approximation。附录把坏事件
\[
\rho(h^*,\hat h_\lambda)\ge\varepsilon
\]
分成风险 shell：
\[
A_l
=
\left\{
h:
2^l\varepsilon^2
\le
\rho^2(h^*,h)
<
2^{l+1}\varepsilon^2
\right\}.
\]

如果 \(\hat h_\lambda\in A_l\)，因为它最小化经验 \(\ell_K\)，经验过程必须出现足够大的偏差：
\[
\sup_{h\in A_l}
\nu_n\{\ell_K(h^\#)-\ell_K(h)\}
\gtrsim
\sqrt n\,(2^l-\tfrac12)\varepsilon^2.
\]

然后用 Bernstein moment condition、variance--mean bound 和 bracketing entropy 积分控制这个概率。
最终得到
\[
\Pr\left(
\rho(h^*,\hat h_\lambda)\ge\varepsilon
\right)
\le
4\exp(-c_en\varepsilon^2).
\]

\subsection{这个证明模块的意义}

这部分并不强依赖 score 的特殊几何。它更像一个通用的 \(M\)-estimator 证明模板：
\begin{enumerate}
\item population loss 必须有正交分解；
\item finite-anchor loss 必须是理想 loss 的小扰动；
\item loss class 需要 variance--mean 和 Bernstein 条件；
\item bracketing entropy 积分要能解出 \(\varepsilon\) 的量级。
\end{enumerate}

因此，这个模块是最容易迁移到你的 velocity matching 方案中的部分。

\section{定理三：有限样本生成误差}

定理三把前面的结果合并为：
\[
\E W_2(P_{Y_0},P_{\widetilde Y_0})
\le
C_{\mathrm{PB}}(t)
\left[
c_1(WL)^{-\frac{2(\eta+1)}{d^*}}
+
c_2\sigma_t^{-(\eta+2)}
\left(
\frac{(WL)^2\log^5(WL)}{n}
\right)^{\frac{\eta+1}{2\eta}}
+
c_3\varepsilon(\widetilde\pi,t,K)
\right].
\]

三项分别对应：
\begin{enumerate}
\item \textbf{神经网络逼近误差}：
\[
(WL)^{-\frac{2(\eta+1)}{d^*}}.
\]
这来自 \(h^*\in C^{\eta+1}\) 的 ReLU 逼近率。
\item \textbf{统计估计误差}：
\[
\sigma_t^{-(\eta+2)}
\left(
\frac{(WL)^2\log^5(WL)}{n}
\right)^{\frac{\eta+1}{2\eta}}.
\]
这来自 bracketing entropy 和经验过程控制。
\item \textbf{anchor 近似误差}：
\[
\varepsilon(\widetilde\pi,t,K).
\]
它衡量有限 anchor/self-normalized importance sampling 和理想 posterior score 之间的差距。
\end{enumerate}

选择
\[
WL
\asymp
\left(
\frac{n}{\log^5 n}
\right)^{\frac{d^*}{2(2\eta+d^*)}}
\]
可平衡前两项，得到
\[
\E W_2(P_{Y_0},P_{\widetilde Y_0})
\lesssim
C_{\mathrm{PB}}(t)
\left[
\left(
\frac{n}{\log^5 n}
\right)^{-\frac{\eta+1}{2\eta+d^*}}
+
\varepsilon(\widetilde\pi,t,K)
\right].
\]

这里最重要的是指数只含 intrinsic dimension \(d^*\)，不是 ambient dimension \(D\)。

\section{对你的随机插值一步生成想法的移植性判断}

你上面的想法是使用随机插值速度场，例如
\[
I_\tau
=
(1-\tau)X_0+\tau X_1+\gamma_\tau\varepsilon,
\qquad
\gamma_\tau=\sqrt{2\tau(1-\tau)},
\]
并用一步生成器
\[
X_1^\theta=h_\theta(U)
\]
诱导模型插值
\[
I_{\theta,\tau}
=
(1-\tau)X_0+\tau h_\theta(U)+\gamma_\tau\varepsilon.
\]

模型速度场为
\[
v_{\theta,\tau}(y)
=
\E\left[
-X_0+h_\theta(U)+\dot\gamma_\tau\varepsilon
\mid I_{\theta,\tau}=y
\right].
\]

\subsection{结论先行}

\textbf{不能原封不动移植第 3 部分的全部证明，但可以移植其中相当大的一部分。}

更精确地说：
\begin{enumerate}
\item 如果你的 ``velocity'' 版本其实使用的是 Gaussian corruption path
\[
Y_\tau=\alpha_\tau X+\sigma_\tau\varepsilon,
\]
那么 velocity matching 与 denoising posterior matching 只差一个确定系数。此时第 3 部分基本可以移植。
\item 如果你使用真正的三项 stochastic interpolant
\[
I_\tau=a_\tau X_0+b_\tau X_1+\gamma_\tau\varepsilon,
\]
那么经验风险/anchor 近似/神经网络逼近这些模块可迁移，但最核心的 single-level pull-back bound 需要重新证明。
\end{enumerate}

\section{为什么 Gaussian velocity 版本几乎可以直接继承}

考虑
\[
Y_\tau=\alpha_\tau X+\sigma_\tau\varepsilon.
\]
sample-wise velocity 是
\[
\dot Y_\tau
=
\dot\alpha_\tau X+\dot\sigma_\tau\varepsilon.
\]
因为
\[
\varepsilon
=
\frac{y-\alpha_\tau X}{\sigma_\tau},
\]
所以真实速度场为
\[
v_{\mathrm{data},\tau}(y)
=
\frac{\dot\sigma_\tau}{\sigma_\tau}y
+
\left(
\dot\alpha_\tau
-
\frac{\alpha_\tau\dot\sigma_\tau}{\sigma_\tau}
\right)
\E[X\mid Y_\tau=y].
\]

模型速度场同理：
\[
v_{\theta,\tau}(y)
=
\frac{\dot\sigma_\tau}{\sigma_\tau}y
+
\left(
\dot\alpha_\tau
-
\frac{\alpha_\tau\dot\sigma_\tau}{\sigma_\tau}
\right)
m_{\theta,\tau}(y).
\]

因此 velocity loss 与 denoiser mismatch 等价：
\[
\mathcal L_{\mathrm{vel}}(\theta)
=
\left(
\dot\alpha_\tau
-
\frac{\alpha_\tau\dot\sigma_\tau}{\sigma_\tau}
\right)^2
\E\|m_{\theta,\tau}(Y_\tau)-X\|^2.
\]

而 denoiser mismatch 又通过 Tweedie 公式等价于 Fisher divergence。于是第 3 部分的 proof chain 仍然成立：
\[
\text{velocity loss}
\Longleftrightarrow
\text{denoiser mismatch}
\Longleftrightarrow
\text{Fisher mismatch}
\Longrightarrow
W_2.
\]

\section{为什么三项 stochastic interpolant 不能直接继承}

对真正的 stochastic interpolant：
\[
I_\tau=a_\tau X_0+b_\tau X_1+\gamma_\tau\varepsilon,
\]
模型密度是
\[
p_{\theta,\tau}(y)
=
\iint
\phi_{\gamma_\tau}
\left(
y-a_\tau x_0-b_\tau h_\theta(u)
\right)
p_0(x_0)\pi(u)\dd x_0\dd u.
\]

速度场是条件期望：
\[
v_{\theta,\tau}(y)
=
\E[
\dot a_\tau X_0+\dot b_\tau h_\theta(U)+\dot\gamma_\tau\varepsilon
\mid I_{\theta,\tau}=y].
\]

这时有一个好消息：velocity matching 仍然有正交分解。若
\[
v_{\mathrm{data},\tau}(y)
=
\E[\dot I_\tau\mid I_\tau=y],
\]
则
\[
\E\|v_{\theta,\tau}(I_\tau)-\dot I_\tau\|^2
=
\E\|v_{\mathrm{data},\tau}(I_\tau)-\dot I_\tau\|^2
+
\E\|v_{\theta,\tau}(I_\tau)-v_{\mathrm{data},\tau}(I_\tau)\|^2.
\]

所以第 3 部分的经验风险证明思路可以迁移。

但坏消息是：这个 excess risk 不再天然等于 Fisher divergence。也就是说，一般没有
\[
\E\|v_{\theta,\tau}-v_{\mathrm{data},\tau}\|^2
\asymp
\mathcal J(p_\tau\|p_{\theta,\tau}).
\]
因此论文中用到的
\[
\text{LSI}+\text{Fisher}
\Longrightarrow
W_2(p_\tau,p_{\theta,\tau})
\]
这一环不能直接使用。

\section{如果要移植，需要补的新定理}

三项 stochastic interpolant 方案要获得 MAGT 第 3 节同等级的理论，至少需要补一个新的 single-level velocity pull-back theorem。

理想形式可能是：
\[
W_2(P_{X_1},P_{h_\theta(U)})
\le
A(\tau)W_2(p_\tau,p_{\theta,\tau})
+
B(\tau)
\|v_{\mathrm{data},\tau}-v_{\theta,\tau}\|_{L^2(p_\tau)}.
\]

但要从固定单层 \(\tau\) 推回 endpoint \(\tau=1\)，还需要一个类似 score-gap growth lemma 的 velocity-gap growth lemma：
\[
\frac{\dd}{\dd s}
\|v_s-v_s^\theta\|_{L^2(p_s)}^2
\le
C(s)
\|v_s-v_s^\theta\|_{L^2(p_s)}^2.
\]

这一步在 MAGT 中由 VP probability flow 的 score PDE 和 Hessian bound 完成；在一般 stochastic interpolant 中并没有现成结论，需要重新建立。

\section{维度率能否保留}

这是最容易被忽略的地方。MAGT 的 rate 依赖 intrinsic dimension \(d^*\)，因为模型平滑的是
\[
\alpha_t h(U)+\sigma_tZ,
\]
低维结构来自 \(h(U)\)，Gaussian noise 只是为了让 ambient density 存在。

而三项 stochastic interpolant 有
\[
I_\tau=(1-\tau)X_0+\tau h(U)+\gamma_\tau\varepsilon.
\]
如果 \(X_0\sim\cN(0,I_D)\) 是全维 Gaussian，那么中间分布天然含有全维 base component。这样一来，证明中的复杂度和几何常数很可能重新依赖 ambient dimension \(D\)，从而失去 MAGT 最重要的 intrinsic-dimension rate。

要尽量保留 intrinsic-dimension 优势，可以考虑：
\begin{enumerate}
\item 让 \(X_0\) 也是低维生成的，例如 \(X_0=g_0(V)\)，而不是全维 Gaussian；
\item 令 \(X_0\) 与 \(U\) 共享低维 latent；
\item 选择靠近终点的固定 \(\tau\)，使 \((1-\tau)X_0\) 足够小；
\item 使用 Gaussian-corruption velocity 版本，而不是完整三项 bridge；
\item 证明新的 effective dimension bound，说明 \(X_0\) 项不会主导复杂度。
\end{enumerate}

\section{可移植模块清单}

\begin{center}
\begin{tabular}{p{0.33\linewidth}p{0.23\linewidth}p{0.34\linewidth}}
\toprule
证明模块 & 可移植性 & 说明 \\
\midrule
posterior anchor 近似 & 高 & 三项 interpolant 可用 \((x_{0,k},u_k)\) anchors。\\
经验风险正交分解 & 高 & velocity matching 也有 conditional expectation projection。\\
bracketing entropy 经验过程 & 中高 & 损失类、Bernstein、variance--mean 需按 velocity loss 重写。\\
ReLU 逼近 \(h^*\) & 高 & 只要仍学习 endpoint map \(h_\theta\)。\\
LSI 推送 & 中 & 若 \(X_0\) 全维，LSI 易得但 intrinsic dimension 可能丢失。\\
Hessian/tube 几何引理 & 中低 & Gaussian corruption endpoint 可用；三项 bridge 需重做。\\
single-level pull-back & 低 & 核心步骤，三项 velocity 方案必须新证。\\
intrinsic-dimension \(W_2\) rate & 不确定 & 取决于 \(X_0\) 的维度和新 pull-back 定理。\\
\bottomrule
\end{tabular}
\end{center}

\section{建议的理论路线}

如果你要把这个 idea 写成可证明的方法，建议分两条路线。

\subsection{路线 A：保守可证路线}

使用 Gaussian-corruption velocity：
\[
Y_\tau=\alpha_\tau X+\sigma_\tau\varepsilon.
\]
把 velocity target 写成 denoiser posterior mean 的线性变换。这样可以最大限度继承 MAGT 第 3 节：
\[
\mathcal L_{\mathrm{vel}}
\asymp
\E\|m_{\theta,\tau}(Y_\tau)-X\|^2
\asymp
\mathcal J(p_\tau\|p_{\theta,\tau}).
\]

这条路线理论最稳，但 stochastic interpolant 的新意较弱。

\subsection{路线 B：真正 stochastic interpolant 路线}

使用
\[
I_\tau=(1-\tau)X_0+\tau X_1+\sqrt{2\tau(1-\tau)}\varepsilon.
\]
证明顺序应当是：
\begin{enumerate}
\item 建立 anchor velocity estimator 的一致性；
\item 建立 velocity loss 的正交分解；
\item 建立 finite-anchor perturbation bound；
\item 建立 velocity-loss class 的 bracketing entropy；
\item 证明新的 single-level velocity pull-back；
\item 分析 \(X_0\) 项对 intrinsic dimension rate 的影响。
\end{enumerate}

这条路线更有创新性，但证明工作量明显更大。

\section{最终判断}

对你上面的想法，第 3 部分证明的移植结论是：
\begin{center}
\fbox{
\begin{minipage}{0.88\linewidth}
\centering
可部分移植；若使用 Gaussian velocity 版本基本可移植；若使用三项 stochastic interpolant，则核心 single-level pull-back 需要重证。
\end{minipage}
}
\end{center}

如果目标是尽快形成一个严谨、可写论文的方法，建议先走 Gaussian fixed-level velocity 版本，作为 MAGT 的 velocity-field reinterpretation；然后再把三项 stochastic interpolant 作为扩展，重点证明新的 single-level velocity pull-back theorem。

\end{document}
